%%-*-latex-*-

\documentclass[wide]{slides}

% Language
%
\usepackage{hyphenat}          % \hyp{} is a breakable dash
\usepackage{alltt}

\frenchspacing  % Follow French conventions after a period

% Maths & Logic
%
\usepackage{stmaryrd}
\usepackage{mathpartir}

% Graphics
%
\usepackage{graphicx}

% Code listing
%
\usepackage{listings}

\usepackage{color}
\definecolor{lightgray}{rgb}{.9,.9,.9}
\definecolor{darkgray}{rgb}{.4,.4,.4}
\definecolor{purple}{rgb}{0.65, 0.12, 0.82}

\lstdefinelanguage{JavaScript}{
  keywords={const, let, typeof, new, true, false, catch, function, return, null, catch, switch, var, if, in, while, do, else, case, break},
  keywordstyle=\color{blue}\bfseries,
  ndkeywords={class, export, boolean, throw, implements, import, this},
  ndkeywordstyle=\color{darkgray}\bfseries,
  identifierstyle=\color{black},
  sensitive=false,
  comment=[l]{//},
  morecomment=[s]{/*}{*/},
  commentstyle=\color{purple}\ttfamily,
  stringstyle=\color{red}\ttfamily,
  morestring=[b]',
  morestring=[b]"
}

\lstset{
   language=JavaScript,
   backgroundcolor=\color{lightgray},
   extendedchars=true,
   basicstyle=\footnotesize\ttfamily,
   showstringspaces=false,
   showspaces=false,
   numbers=left,
   numberstyle=\footnotesize,
   numbersep=9pt,
   tabsize=2,
   breaklines=true,
   showtabs=false,
   captionpos=b
}

% New environments and commands
%

% Input files
%
% \input{commands}

% ----------------------------------------------------------------
% Document
%
\maintitle{LIGO for TypeScript developers on Layer-1}
\mainauthor{\textbf{Christian Rinderknecht}\\
  {\small\url{Christian.Rinderknecht@trili.tech}}\\
\textsf{TriliTech}}
\confname{Show and Tell}
\confshortname{S\&T}
\confdate{17 April 2025}

\begin{document}

\maketitle

\begin{slide}
  \title{A Short History of LIGO}

  \begin{itemize}

    \item I did my PhD at INRIA under the supervision of Michel Mauny.

    \item I joined Nomadic in July 2018, where I started designing a
      smart contract language, similar to \textsf{Pascal}, that would
      be compiled to Michelson.

    \item My then-colleague Gabriel Alfour also wanted to make
      programming smart contracts on \textsf{Tezos} easier, but with a
      bottom\hyp{}up approach: building a set of macros that expand
      into Michelson.

    \item We joined our efforts in 2019 into what became LIGO.

    \item I contributed the design and the front-end, he did the
      type-checker and code generator (50-50 in lines of code).

  \end{itemize}
\end{slide}

\begin{slide}
  \title{A Short History of LIGO (continued)}

  \begin{itemize}

    \item The idea was for LIGO to offer different concrete syntaxes
      inspired from mainstream programming languages.

    \item Those LIGO languages were meant to be domain-specific, but
      not overly so: supporting special literals (natural numbers,
      mutez, for examples), but taking most Tezos-specific types and
      functions from a library.

    \item The front-end had to lower their input to the common
      pipeline, while sharing as much code as possible between the
      processing of the different syntaxes.

    \item The first syntax was \textsf{PascaLIGO}, as I thought it
      would suit an imperative style, and lots of keywords would make
      the code more explicit, given the high stakes.

    \item Next, turmoil in the nascent ecosystem lead me to design
      \textsf{CameLIGO} (internally called \textsf{Ligodity}). Then
      came \textsf{ReasonLIGO}, \textsf{JsLIGO} (which is today's
      topic), and even \textsf{PyLIGO}, which was never plugged in.

  \end{itemize}
\end{slide}

\begin{slide}
  \title{Architecture}

  \begin{itemize}

    \item Nowadays, only two LIGO syntaxes remain: \textsf{CameLIGO}
      and \textsf{JsLIGO}. The former is a strict subset of
      \textsf{OCaml}, if we leave aside the Tezos-specific literals, a
      local type definition construct, and attributes (they occur in
      prefix position in LIGO).

    \item \textsf{JsLIGO} started also as a subset but evolved and
      became syntactically incompatible with its model,
      \textsf{TypeScript}. For example, a proposal for pattern
      matching in \textsf{TypeScript} was implemented in
      \textsf{JsLIGO}.

    \item The front-end is made, in order, of a preprocessor, a lexer,
      two parsers and the first tree transformation before
      type\hyp{}inference. Each pass can be built as a standalone
      executable for testing.

    \item I wrote the preprocessor because, at the time, we had no
      time to design and implement a module system, together with a
      build system.

    \item Unfortunately, it became very popular, even when LIGO got
      itself a module system, because it enabled conditional
      compilation (often for new type definitions of the storage,
      depending on the version of the standard being implemented).

  \end{itemize}
\end{slide}

\begin{slide}
  \title{Preprocessor and Lexer}

  \begin{itemize}

    \item The preprocessor is not trivial, as it needs to understand
      UTF-8 encoded glyphs and the kind of comments and strings that
      depend on the LIGO variant.

    \item The lexer is also a functor parameterised by lexers that are
      specific to a given LIGO syntax (comments, strings, keywords
      etc.). Error reporting has to be aware of linemarkers left by
      the preprocessor. Some style constraints are enforced (like a
      linter) for the sake of clarity.

    \item The lexer is specified as an input to \textsf{ocamllex}.

    \item Until \textsf{JsLIGO} was added to the party, the lexer
      signature was a function that returned a single token, which the
      parser called on demand. Syntactic constructs difficult to parse
      sometimes rely on lexer hacks, that is, having the parser
      informing the lexer of the syntactic left\hyp{}context. Instead,
      I had the lexer run on all the (preprocessed) input, before
      applying a series of self\hyp{}passes on the tokens
      ---~injecting virtual tokens when needed to help the
      parser. This is not as expressive as having the syntactic
      context, but bi\hyp{}directional, unbounded lexical context
      works quite well in practice.

  \end{itemize}
\end{slide}

\begin{slide}
  \lstset{language=JavaScript}

  \title{Lexer}

  \begin{itemize}

    \item Perhaps the most interesting self\hyp{}pass on the tokens is
      found with \textsf{JsLIGO}:
      \begin{lstlisting}
        const f = ([x,y]) ...
      \end{lstlisting}
      could be a parenthesised pair or the start of an arrow
      function. Analysing the lexical right\hyp{}context helps
      disambiguate, and the lexer injects a \texttt{ES6FUN} token in
      front of a function. The parser then knows a function is coming
      next upon reading the virtual token \texttt{ES6FUN}.

    \item The LIGO pipeline is more than a compiler, because it
      provides support for a LSP, and that is why the front-end has to
      keep all the source information, including comments. This adds
      significant complexity to the lexer, which has to hide it from
      the parser.

  \end{itemize}
\end{slide}

\begin{slide}
  \title{Parser}

  \begin{itemize}

    \item The two LIGO parsers are generated by \textsf{Menhir}.

    \item \textsf{Menhir} accepts LR(1) grammars, with semantic
      actions in \textsf{OCaml}. In practice, many grammars are
      annotated with precedence declarations to solve
      LR\hyp{}conflicts, potentially making the grammars not LR(1),
      and making it difficult to reason about them.

    \item Both \textsf{CameLIGO} and \textsf{JsLIGO} are given pure
      LR(1) grammars, thanks to the rewriting of productions, and the
      use of virtual tokens.

    \item \textsf{Menhir} can identify all error states in the
      LR\hyp{}automaton, so we can write precise error messages for
      them. They are precise because the syntactic context is
      available at the error state, instead of only lexical context,
      so the past (the intention of the programmer) can be expressed
      in terms of constructs, and the future too if we take some care
      when writing the grammar.

  \end{itemize}
\end{slide}

\begin{slide}
  \title{Parser and Unified AST}

  \begin{itemize}

    \item The parser produces an Abstract Syntax Tree (AST) ---~named
      Concrete Syntax Tree (CST) in the code base for legacy reasons.

    \item The AST is then transformed in a simpler kind of tree, the
      \emph{unified AST}.

    \item Whereas the ASTs are specific to each LIGO variant, the
      unified AST is common to all, and can be considered as the
      entrance to the common pipeline.

  \end{itemize}
\end{slide}



\end{document}
